%% LyX 2.4.4 created this file.  For more info, see https://www.lyx.org/.
%% Do not edit unless you really know what you are doing.
\documentclass[english]{article}
\usepackage[T1]{fontenc}
\usepackage[latin9]{inputenc}
\usepackage{enumitem}
\usepackage{amsmath}
\usepackage{amsthm}
\usepackage{amssymb}
\usepackage{geometry}
\geometry{verbose,tmargin=1in,bmargin=1in,lmargin=1in,rmargin=1in}
\PassOptionsToPackage{normalem}{ulem}
\usepackage{ulem}

\makeatletter
%%%%%%%%%%%%%%%%%%%%%%%%%%%%%% Textclass specific LaTeX commands.
\numberwithin{equation}{section}
\numberwithin{figure}{section}
\newlength{\lyxlabelwidth}      % auxiliary length 
\theoremstyle{plain}
\newtheorem*{thm*}{\protect\theoremname}
\theoremstyle{plain}
\newtheorem*{prop*}{\protect\propositionname}
\theoremstyle{definition}
\newtheorem*{example*}{\protect\examplename}

%%%%%%%%%%%%%%%%%%%%%%%%%%%%%% User specified LaTeX commands.
\usepackage{fancyhdr}
\usepackage{lastpage}
\pagestyle{fancy}
\usepackage{enumitem}
\usepackage{ifthen}
\everymath=\expandafter{\the\everymath\displaystyle}
%\DeclareMathSizes{10}{10}{10}{10}
\usepackage{multicol}

\usepackage{tikz}
\usepackage{tikz-3dplot}
\usetikzlibrary{calc,arrows}

\usetikzlibrary{patterns}
\usetikzlibrary{plotmarks}
\usepackage{pgfplots}
\pgfplotsset{compat=newest}

\usepackage{calc}
\usepackage[nomessages]{fp}

\usepackage{datetime}
\usepackage[super]{nth}
% Added by lyx2lyx
\setlength{\parskip}{\smallskipamount}
\setlength{\parindent}{0pt}

\makeatother

\usepackage{babel}
\providecommand{\examplename}{Example}
\providecommand{\propositionname}{Proposition}
\providecommand{\theoremname}{Theorem}

\begin{document}
\renewcommand{\author}{Dr.\,James Ashe}

\newcommand{\classNum}{336}

\newcommand{\classSec}{A}

\newcommand{\examType}{Exam 2 Review}

\renewcommand{\date}{\formatdate{23}{10}{2013}}

%%First page's header%%%%%%%%%%%%%%%%%%%%%%
\fancypagestyle{firststyle} %%header settings for first pagr
{
	\fancyhead[L]{MTH \classNum{}(\classSec{}) \author{}\\
	\textbf{\examType{}} ~\date{}}
	\fancyhead[C]{}
	\fancyhead[R]{\textbf{}}
	\setlength{\headsep}{14pt}
}
%%%%%%%%%%%%%%%%%%%%%%%%%%%%%%%%%%%%%%%%%%

%%Next page's header%%%%%%%%%%%%%%%%%%%%%%%
\setlist{nolistsep}
\lhead{\classNum{}(\classSec{})} 
\chead{\textbf{}} 
\cfoot{\thepage{}~of~\pageref{LastPage}} 
\setlength{\headsep}{5pt} 
%%%%%%%%%%%%%%%%%%%%%%%%%%%%%%%%%%%%%%%%%%%

%%Various settings; see the preamble for other settings%%%%%
%\setenumerate[0]{leftmargin=12pt,itemindent=0pt,labelwidth=0pt}
\setlist{topsep=.5pt}
\renewcommand{\labelenumi}{\textbf{\arabic{enumi}.}}
\renewcommand{\labelenumii}{\textbf{\alph{enumii}.}}
\thispagestyle{firststyle}
%%%%%%%%%%%%%%%%%%%%%%%%%%%%%%%%%%%%%%%%%%

\newcommand{\xmax}{0}
\newcommand{\xmin}{-\xmax}
\newcommand{\xstep}{0}
\newcommand{\ymax}{0}
\newcommand{\ymin}{-\ymax}
\newcommand{\ystep}{0} 
\newcommand{\dspace}{\vspace{1cm}}
\newcommand{\xa}{0}
\newcommand{\xb}{0}
\newcommand{\ya}{1}
\newcommand{\yb}{1}

\subsubsection*{Chapter 2 Matrix Algebra}
\begin{itemize}
\item \textbf{2.1 Matrix operations: }be able to compute the basic operations
and know when the operation is undefined. 
\item \textbf{2.2 Linear transformations: }

\begin{itemize}
\item Find the standard matrix for a given a given linear transformation,
e.g., HW 2.

\begin{itemize}
\item Find the columns of the standard matrix by using the column vectors
from the standard basis., e.g., $T\left(\begin{bmatrix}\begin{array}{r}
1\\
0\\
0
\end{array}\end{bmatrix}\right)$, $T\left(\begin{bmatrix}\begin{array}{r}
0\\
1\\
0
\end{array}\end{bmatrix}\right)$, and $T\left(\begin{bmatrix}\begin{array}{r}
0\\
0\\
1
\end{array}\end{bmatrix}\right)$ gives the \nth{1}, \nth{2}, and \nth{3} columns respectively.
\end{itemize}
\item Determine if a given function is a \emph{linear transformation}, e.g.,
HW 4ab.

\begin{itemize}
\item Recall that to prove that $T$ is a linear transformation it suffices
to find a standard matrix, $A\in\mathcal{M}_{m\times n}$ for $\mu_{A}:\mathbb{R}^{n}\rightarrow\mathbb{R}^{m}$.
To prove that it is not a standard matrix you need to show that it
fails at least one of the properties.
\end{itemize}
\end{itemize}
\item \textbf{2.3 Inverse matrices. }

\begin{itemize}
\item Use Gaussian elimination to determine if a matrix $A$ is invertible
and find $A^{-1}$.
\item Use the inverse $A^{-1}$ to solve $A\mathbf{x}=\mathbf{b}$ e.g.,
HW 2ab.
\end{itemize}
\begin{thm*}
An $n\times n$ matrix is \uline{nonsingular}\emph{ $\iff$ }it is
invertible. 
\end{thm*}
\end{itemize}

\subsubsection*{Chapter 3 Vector Spaces}
\begin{itemize}
\item \textbf{3.1 Subspaces of $\mathbb{R}^{n}$.}

\begin{itemize}
\item Determine if a given set is a \emph{subspace}, e.g. HW 1ab.
\end{itemize}
\begin{prop*}
Let $A\in\mathcal{M}_{m\times n}$. The set of solutions to $A\mathbf{x}=\mathbf{0}$
is a subspace of $\mathbb{R}^{n}$
\end{prop*}
\vspace{0cm}
\begin{prop*}
Let $\mathbf{v}_{1},\mathbf{v}_{2},\dots,\mathbf{v}_{k}\in\mathbb{R}^{n}$.
Then $V=\mbox{Span\ensuremath{\left(\mathbf{v}_{1},\mathbf{v}_{2},\dots,\mathbf{v}_{k}\right)}}$
is a subspace of $\mathbb{R}^{n}$.
\end{prop*}
\begin{itemize}
\item Decide whether a given set of vectors \emph{spans} $\mathbb{R}^{n}$,
e.g., HW 2ab.

\begin{itemize}
\item Note: you can use dimensions (section 3.4) to determine that some
sets do not span $\mathbb{R}^{n}$.
\end{itemize}
\item [$\star$]Let $\mathbf{v}_{1},\mathbf{v}_{2},\dots,\mathbf{v}_{k}\in\mathbb{R}^{n}$
and let $\mathbf{v}\in\mathbb{R}^{n}$. Prove that $\mathbf{v}\in\mbox{Span}\left(\mathbf{v}_{1},\mathbf{v}_{2},\dots,\mathbf{v}_{k}\right)$
$\Rightarrow$ $\mbox{Span}\left(\mathbf{v}_{1},\mathbf{v}_{2},\dots,\mathbf{v}_{k}\right)=\mbox{Span}\left(\mathbf{v}_{1},\mathbf{v}_{2},\dots,\mathbf{v}_{k},\mathbf{v}\right)$.
\end{itemize}
\item \textbf{3.3 Linear Independence and Basis}

\begin{prop*}
\textup{\emph{Let $A\in\mathcal{M}_{n\times n}$. Then $A$ is nonsingular
if and only if its column vectors form a basis for $\mathbb{R}^{n}$.}}
\end{prop*}
\begin{itemize}
\item Determine if a given set of vectors is \emph{linearly independent},
e.g., HW 2ab.

\begin{itemize}
\item Find the reduced echelon form of the matrix of column vectors; $A\mathbf{x}=\mathbf{0}$
has only the trivial solution$\Rightarrow$the vectors are linearly
independent.
\end{itemize}
\item Determine if a given set of vectors is a \emph{basis }for a given
space, e.g., HW 3bc

\begin{itemize}
\item To show that $\mathbf{v}_{1},\mathbf{v}_{2},\dots,\mathbf{v}_{k}\in\mathbb{R}^{n}$
is a basis, you must show that $\left\{ \mathbf{v}_{1},\mathbf{v}_{2},\dots,\mathbf{v}_{k}\right\} $
are linearly independent (see above) \uline{and} $\mbox{Span}\{\mathbf{v}_{1},\mathbf{v}_{2},\dots,\mathbf{v}_{k}\}=\mathbb{R}^{n}$.
If the set fails \uline{either} of these, it is not a basis for $\mathbb{R}^{n}$. 
\item Section 3.4 can be used to show that a set is not a basis.\newpage
\end{itemize}
\end{itemize}
\item \textbf{3.4 Dimension and Its Consequences. }

\begin{prop*}
Let $V\subset\mathbb{R}^{n}$ be a subspace, let $\{\mathbf{v}_{1},\dots,\mathbf{v}_{k}\}$
be a basis for $V$, and let $\{\mathbf{w}_{1},\dots,\mathbf{w}_{\ell}\}\in V$.
If $\ell>k$ then $\{\mathbf{w}_{1},\dots,\mathbf{w}_{\ell}\}$ must
be \uline{linearly dependent}.

\begin{itemize}
\item $\dim\mathbb{R}^{n}=n$\emph{; meaning that a basis for $\mathbb{R}^{n}$
such as the standard basis will have $n$ vectors.}
\item \emph{Any two bases for a space will have the same number of vectors. }
\item \emph{Let }$\{\mathbf{v}_{1},\dots,\mathbf{v}_{k}\}\subset\mathbb{R}^{n}$\emph{.
If $k\neq n$ then }$\{\mathbf{v}_{1},\dots,\mathbf{v}_{k}\}$ \emph{is
}\emph{\uline{not}}\emph{ a basis for $\mathbb{R}^{n}$. If $k=n$
then it is a basis if and only if }$\{\mathbf{v}_{1},\dots,\mathbf{v}_{k}\}$
\emph{linearly independent. }
\end{itemize}
\begin{example*}
(HW 1.b) Find a basis for each of the given subspaces and determine
its dimension.
\end{example*}
\begin{enumerate}
\item $V=\left\{ \mathbf{x}\in\mathbb{R}^{n}\,:\,x_{1}+x_{2}+x_{3}+x_{4}=0,\,x_{2}+x_{4}=0\right\} \subset\mathbb{R}^{4}$.
\[
\begin{bmatrix}\begin{array}{rrrr|r}
1 & 1 & 1 & 1 & 0\\
0 & 1 & 0 & 1 & 0
\end{array}\end{bmatrix}\rightsquigarrow\begin{bmatrix}\begin{array}{rrrr|r}
1 & 0 & 1 & 0 & 0\\
0 & 1 & 0 & 1 & 0
\end{array}\end{bmatrix}\Rightarrow x_{1}=-x_{3}\mbox{ and }x_{2}=-x_{4}
\]
$\Rightarrow$the general solution is $\mathbf{x}=x_{3}(-1,0,1,0)+x_{4}(0,-1,0,1)$.
So $V$ is spanned by $(-1,0,1,0)$ and \textup{$(0,-1,0,1)$. It's
easily checked that these vectors are linearly independent. So $V$
is a basis with $\dim V=2$.}
\item $V=\mbox{span}\left\{ (3,-2,4),(-3,2,-4),(-6,4,-8)\right\} $. 
\[
\begin{bmatrix}\begin{array}{rrr}
3 & -3 & -6\\
-2 & 2 & 4\\
4 & -4 & -8
\end{array}\end{bmatrix}\rightsquigarrow\begin{bmatrix}\begin{array}{rrr}
1 & -1 & -2\\
0 & 0 & 0\\
0 & 0 & 0
\end{array}\end{bmatrix}
\]
So $(3,-2,4)$ (as is the other vectors) is a basis and $\dim V=1$.
\end{enumerate}
\end{prop*}
\end{itemize}

\end{document}
